\Askhsh{22098}{4}
\begin{ylh}{1}
    \item 9.2. Το Πυθαγόρειο θεώρημα
    \item 9.4. Γενίκευση του Πυθαγόρειου θεωρήματος
    \item 10.3. Εμβαδόν βασικών ευθύγραμμων σχημάτων
    \item 11.7. Εμβαδόν κυκλικού τομέα και κυκλικού τμήματος.
\end{ylh}
\begin{wrapfigure}[9]{r}{5.5cm} % The [9] can be adjusted if the figure takes up more/less lines.
\centering
\begin{tikzpicture}[scale=0.75]
    % Define a scaling factor for 'a' to draw the figure with specific dimensions.
    % We set \alpha = 1 for the drawing coordinates.
    \pgfmathsetmacro{\drawingAlpha}{1} 

    % Define points for the rectangle ABΓΔ based on the dimensions AB = 4\alpha, AΔ = \pi\alpha.
    % We place Δ at the origin (0,0).
    \tkzDefPoint(0, 0){D} % Point Δ (Delta)
    \tkzDefPoint({4*\drawingAlpha}, 0){C} % Point Γ (Gamma)
    \tkzDefPoint({4*\drawingAlpha}, {pi*\drawingAlpha}){B} % Point B
    \tkzDefPoint(0, {pi*\drawingAlpha}){A} % Point A

    % Define M as the midpoint of AB (the center of the semicircle).
    \tkzDefMidPoint(A,B) \tkzGetPoint{M}

    % Draw the rectangle ABΓΔ.
    \tkzDrawPolygon(A,B,C,D)

    % Draw the semicircle with diameter AB.
    % The problem states it's "inside the rectangle", so the arc should be drawn downwards from AB.
    % The center of the semicircle is M, and the radius is AB/2 = 2\alpha.
    % We draw from B (angle 0 relative to M) to A (angle -180 relative to M) in a clockwise direction.
    \draw (M) arc (0:-180:{2*\drawingAlpha});

    % Draw points A, B, M, Δ, Γ.
    \tkzDrawPoints(A,B,M,D,C)

    % Label the points.
    \tkzLabelPoint[above left](A){A}
    \tkzLabelPoint[above right](B){B}
    \tkzLabelPoint[above](M){M}
    \tkzLabelPoint[below left](D){$\Delta$}
    \tkzLabelPoint[below right](C){$\Gamma$}
\end{tikzpicture}
\caption{Το ορθογώνιο ABΓΔ με το ημικύκλιο.}
\end{wrapfigure}
Στο παρακάτω σχήμα το ΑΒΓΔ είναι ορθογώνιο με $\text{AB} = 4\alpha$ και $\text{AΔ} = \pi\alpha$. Στο εσωτερικό του ορθογωνίου σχεδιάστηκε ημικύκλιο διαμέτρου ΑΒ.
\begin{alist}
    \item Να αποδείξετε ότι το ημικύκλιο χωρίζει το ορθογώνιο σε δύο ισεμβαδικά χωρία. \monades{8}
    \item Αν η διαγώνιος ΒΔ τέμνει το ημικύκλιο στο σημείο Ε και Μ είναι το μέσο της ΑΒ,
    \begin{rlist}
        \item να αποδείξετε ότι $\text{AB}^2 = \text{BΔ} \cdot \text{BE}$ και $\text{AΔ}^2 = \text{BΔ} \cdot \text{ΔE}$. \monades{6}
        \item να αποδείξετε ότι $\text{BE} = \dfrac{16\alpha}{\sqrt{16+\pi^2}}$ και $\text{ΔE} = \dfrac{\pi^2\alpha}{\sqrt{16+\pi^2}}$. \monades{6}
        \item να υπολογίσετε το $\text{συν}\widehat{\text{BME}}$. \monades{5}
    \end{rlist}
\end{alist}